\par\Needspace{10\baselineskip}
\originalchapter{官方作业二}
本章翻译18.950 Homework 2的两个原题, 共20分. 截止日未载于原PDF; 解答为AI独立编写与复审.

\begin{exercise}\label{ps:2-1}
原题PS2.1, 5分. 设正则曲线 $c:I\to\R^2$ 处处满足 $|\kappa_c|\le1$. 沿一个方向拉伸, 定义 $d(t)=(2c_1(t),c_2(t))$. 对 $\kappa_d$ 可以作怎样的判断? 圆是首先应考虑的例子.
\end{exercise}
\begin{solution}\label{pssol:2-1}
设 $A=\operatorname{diag}(2,1)$. 对每个 $t$, $|c'|\le|Ac'|\le2|c'|$, 且 $\det(A)=2$. 所以
\[
 \kappa_d=2\left(\frac{|c'|}{|Ac'|}\right)^3\kappa_c,
 \qquad \frac14|\kappa_c|\le|\kappa_d|\le2|\kappa_c|\le2.
\]
两条曲线的有向曲率同号, 零曲率点也相同. 因而并不能保证拉伸后仍有 $|\kappa_d|\le1$, 但总能保证上界2.

取单位圆 $c(t)=(\cos t,\sin t)$, 得 $d(t)=(2\cos t,\sin t)$. 第1章椭圆公式给出在 $(\pm2,0)$ 处曲率2, 在 $(0,\pm1)$ 处曲率 $1/4$, 同时实现上述两个比例界. 若只知道原曲率的上界, 对新曲率的绝对值并没有正的统一下界, 因为直线的曲率为0.
\end{solution}

\begin{exercise}\label{ps:2-2}
原题PS2.2, 15分. 本题需要复数的初等四则运算, 原题说明还需要MAPLE或Mathematica, 且不要求证明实验观察. 设 $p$ 是多项式, 例如
\[
\begin{aligned}
 p_1(z)&=z^4/4-8z^3/3+17z^2/2-10z,\\
 p_2(z)&=z^4-z,\qquad
 p_3(z)=z^4/4+z^3/3+z^2/2+z.
\end{aligned}
\]
固定半径 $R$, 考察闭合的 $2\pi$ 周期曲线 $c(t)=p(Re^{it})$. 总曲率如何随 $R$ 变化? 跳跃发生在哪里, 有几次? 用计算机画一些图.
\end{exercise}
\begin{solution}\label{pssol:2-2}
原讲义 Definition 8.3 的总曲率是有向积分 $\int\kappa|c'|\dd t$, 而不是 $\int|\kappa||c'|\dd t$. 设 $R>0$, 且 $p$ 非常数. 导数 $c'(t)=iRe^{it}p'(Re^{it})$ 表明, 曲线正则当且仅当半径 $R$ 的圆上没有 $p'$ 的零点. 在这些半径, 第2章的旋转数公式给出
\[
 \int_0^{2\pi}\kappa(t)|c'(t)|\dd t=2\pi\bigl(1+N(R)\bigr),
 \qquad N(R)=\#\{p'\text{的零点 }a:|a|<R\},
\]
其中零点按重数计. 这里附上初等依据: 把 $p'$ 分解为常数与若干 $z-a$ 的乘积. 速度中的 $e^{it}$ 转一周贡献1; 每个因子 $Re^{it}-a$ 在 $|a|<R$ 时绕0一次, 在 $|a|>R$ 时绕0零次. 后一判断可分别把 $a$ 连续缩至0或把圆半径连续缩至0, 过程中因子均不经过0. 乘积的辐角增量等于各因子的增量之和.

对三组指定多项式,
\[
 p_1'(z)=(z-1)(z-2)(z-5),\quad
 p_2'(z)=4z^3-1,\quad p_3'(z)=(z+1)(z^2+1).
\]
于是 $p_1$ 的总有向曲率在区间 $(0,1),(1,2),(2,5),(5,\infty)$ 上依次为 $2\pi,4\pi,6\pi,8\pi$; 在临界半径1、2、5处发生三次跳跃, 每次两侧之差为 $2\pi$. $p_2'$ 的三个单零点均有模 $\rho=4^{-1/3}$, 所以在 $(0,\rho)$ 和 $(\rho,\infty)$ 上分别为 $2\pi,8\pi$, 只有一个跳跃半径, 跳幅 $6\pi$. $p_3'$ 的三个单零点 $-1,i,-i$ 均有模1, 也只有一次跳跃, 两侧同样为 $2\pi$ 和 $8\pi$.

在这些临界半径, $c'$ 在相应参数处为0, 常规正则曲线总曲率公式本身不适用; 上述跳跃指两侧正则半径的值, 不是为奇点半径强行赋值. $R=0$ 也是非正则的常值曲线. 若改问绝对总曲率, 则上面的整数阶梯结论不能直接使用.

下面的计算机绘图由TikZ/PGFPlots按 $p_j(Re^{it})$ 的实部和虚部独立采样生成. 它们完成原题的画图要求, 不冒称使用了MAPLE或Mathematica, 也不以采样图代替零点计数证明. 每幅图的横、纵坐标分别为复数像的实部和虚部, 同一幅保持相等的坐标尺度.
\end{solution}
\par\Needspace{11\baselineskip}
\begin{center}
\begin{tikzpicture}
\begin{axis}[width=6.4cm,height=4.7cm,axis equal image,axis lines=box,xtick=\empty,ytick=\empty,title={$p_1,\ R=1/2$}]
\addplot[thick,structurecolor,domain=0:360,samples=241] ({(0.015625)*cos(4*x)+(-0.333333333333)*cos(3*x)+(2.125)*cos(2*x)+(-5)*cos(1*x)},{(0.015625)*sin(4*x)+(-0.333333333333)*sin(3*x)+(2.125)*sin(2*x)+(-5)*sin(1*x)});
\end{axis}
\end{tikzpicture}
\hspace{4mm}
\begin{tikzpicture}
\begin{axis}[width=6.4cm,height=4.7cm,axis equal image,axis lines=box,xtick=\empty,ytick=\empty,title={$p_1,\ R=3/2$}]
\addplot[thick,structurecolor,domain=0:360,samples=241] ({(1.265625)*cos(4*x)+(-9)*cos(3*x)+(19.125)*cos(2*x)+(-15)*cos(1*x)},{(1.265625)*sin(4*x)+(-9)*sin(3*x)+(19.125)*sin(2*x)+(-15)*sin(1*x)});
\end{axis}
\end{tikzpicture}
\end{center}
\par\Needspace{11\baselineskip}
\begin{center}
\begin{tikzpicture}
\begin{axis}[width=6.4cm,height=4.7cm,axis equal image,axis lines=box,xtick=\empty,ytick=\empty,title={$p_1,\ R=3$}]
\addplot[thick,structurecolor,domain=0:360,samples=241] ({(20.25)*cos(4*x)+(-72)*cos(3*x)+(76.5)*cos(2*x)+(-30)*cos(1*x)},{(20.25)*sin(4*x)+(-72)*sin(3*x)+(76.5)*sin(2*x)+(-30)*sin(1*x)});
\end{axis}
\end{tikzpicture}
\hspace{4mm}
\begin{tikzpicture}
\begin{axis}[width=6.4cm,height=4.7cm,axis equal image,axis lines=box,xtick=\empty,ytick=\empty,title={$p_1,\ R=6$}]
\addplot[thick,structurecolor,domain=0:360,samples=241] ({(324)*cos(4*x)+(-576)*cos(3*x)+(306)*cos(2*x)+(-60)*cos(1*x)},{(324)*sin(4*x)+(-576)*sin(3*x)+(306)*sin(2*x)+(-60)*sin(1*x)});
\end{axis}
\end{tikzpicture}
\end{center}
\par\Needspace{11\baselineskip}
\begin{center}
\begin{tikzpicture}
\begin{axis}[width=6.4cm,height=4.7cm,axis equal image,axis lines=box,xtick=\empty,ytick=\empty,title={$p_2,\ R=2/5$}]
\addplot[thick,structurecolor,domain=0:360,samples=241] ({(0.0256)*cos(4*x)+(-0.4)*cos(1*x)},{(0.0256)*sin(4*x)+(-0.4)*sin(1*x)});
\end{axis}
\end{tikzpicture}
\hspace{4mm}
\begin{tikzpicture}
\begin{axis}[width=6.4cm,height=4.7cm,axis equal image,axis lines=box,xtick=\empty,ytick=\empty,title={$p_2,\ R=9/10$}]
\addplot[thick,structurecolor,domain=0:360,samples=241] ({(0.6561)*cos(4*x)+(-0.9)*cos(1*x)},{(0.6561)*sin(4*x)+(-0.9)*sin(1*x)});
\end{axis}
\end{tikzpicture}
\end{center}
\par\Needspace{11\baselineskip}
\begin{center}
\begin{tikzpicture}
\begin{axis}[width=6.4cm,height=4.7cm,axis equal image,axis lines=box,xtick=\empty,ytick=\empty,title={$p_3,\ R=1/2$}]
\addplot[thick,structurecolor,domain=0:360,samples=241] ({(0.015625)*cos(4*x)+(0.0416666666667)*cos(3*x)+(0.125)*cos(2*x)+(0.5)*cos(1*x)},{(0.015625)*sin(4*x)+(0.0416666666667)*sin(3*x)+(0.125)*sin(2*x)+(0.5)*sin(1*x)});
\end{axis}
\end{tikzpicture}
\hspace{4mm}
\begin{tikzpicture}
\begin{axis}[width=6.4cm,height=4.7cm,axis equal image,axis lines=box,xtick=\empty,ytick=\empty,title={$p_3,\ R=3/2$}]
\addplot[thick,structurecolor,domain=0:360,samples=241] ({(1.265625)*cos(4*x)+(1.125)*cos(3*x)+(1.125)*cos(2*x)+(1.5)*cos(1*x)},{(1.265625)*sin(4*x)+(1.125)*sin(3*x)+(1.125)*sin(2*x)+(1.5)*sin(1*x)});
\end{axis}
\end{tikzpicture}
\end{center}

